\section{The substitution method}

The substitution method is the integration counterpart of the chain rule. It is used when the integrand contains a composite expression together with something close to its derivative.

The idea is to rename the inner expression so that the integral becomes simpler.

\subsection*{The basic idea}

Suppose we see an expression such as
\[
\int 2x\cos(x^2)\,dx.
\]
The inside expression is \(x^2\), and its derivative is \(2x\). This suggests the substitution
\[
u=x^2.
\]
Then
\[
du=2x\,dx.
\]
The integral becomes
\[
\int \cos u\,du=\sin u+C.
\]
Returning to \(x\), we obtain
\[
\int 2x\cos(x^2)\,dx=\sin(x^2)+C.
\]

\begin{remarkbox}
Substitution works by recognizing an inside function and its derivative. It reverses the chain rule.
\end{remarkbox}

\subsection*{A step-by-step example}

\begin{examplebox}
Compute
\[
\int x\sqrt{1+x^2}\,dx.
\]
The expression under the square root is
\[
1+x^2.
\]
Its derivative is \(2x\), and the integrand contains \(x\,dx\). Let
\[
u=1+x^2.
\]
Then
\[
du=2x\,dx,
\qquad
x\,dx=\frac12du.
\]
The integral becomes
\[
\int x\sqrt{1+x^2}\,dx
=\frac12\int u^{1/2}\,du.
\]
Using the power rule,
\[
\frac12\int u^{1/2}\,du
=\frac12\cdot\frac{u^{3/2}}{3/2}+C
=\frac13u^{3/2}+C.
\]
Return to \(x\):
\[
\int x\sqrt{1+x^2}\,dx=\frac13(1+x^2)^{3/2}+C.
\]
\end{examplebox}

The factor \(1/2\) appears because the integrand contains \(x\,dx\), not exactly \(2x\,dx\). Constants can be adjusted.

\subsection*{Substitution with logarithms and exponentials}

Substitution often appears when a logarithm or exponential contains an inner expression.

\begin{examplebox}
Compute
\[
\int \frac{e^x}{2e^x+1}\,dx.
\]
Let
\[
u=2e^x+1.
\]
Then
\[
du=2e^x\,dx,
\qquad
e^x\,dx=\frac12du.
\]
Therefore
\[
\int \frac{e^x}{2e^x+1}\,dx
=\frac12\int \frac1u\,du.
\]
Thus
\[
\int \frac{e^x}{2e^x+1}\,dx
=\frac12\ln|u|+C
=\frac12\ln|2e^x+1|+C.
\]
\end{examplebox}

In this example, the denominator became the new variable. Its derivative was present in the numerator up to a constant factor.

\subsection*{Substitution with trigonometric functions}

\begin{examplebox}
Compute
\[
\int \cos x\,e^{\sin x}\,dx.
\]
The inside expression is \(\sin x\), and its derivative is \(\cos x\). Let
\[
u=\sin x.
\]
Then
\[
du=\cos x\,dx.
\]
The integral becomes
\[
\int e^u\,du=e^u+C.
\]
Returning to \(x\),
\[
\int \cos x\,e^{\sin x}\,dx=e^{\sin x}+C.
\]
\end{examplebox}

The method is the same: identify the inside function, check whether its derivative appears, and rewrite the integral in the new variable.

\subsection*{Substitution is a change of language}

A substitution is not just a symbol trick. It changes the variable so that the structure becomes simpler. The final answer must be returned to the original variable unless the problem asks otherwise.

\begin{summarybox}
When trying substitution, ask:
\begin{itemize}
\item Is there an inside expression repeated in the integrand?
\item Is its derivative present, possibly up to a constant factor?
\item What should be called \(u\)?
\item What is \(du\)?
\item After integrating in \(u\), did I return to the original variable?
\item Can I check by differentiating the final answer?
\end{itemize}
\end{summarybox}

Substitution is one of the most important integration methods because many complicated-looking integrals are simple chain rules read backward.